% Autogenerated translation of LANGUAGE_SPEC_v1.md by Texpad
% To stop this file being overwritten during the typeset process, please move or remove this header

\documentclass[12pt]{book}
\usepackage{graphicx}
\usepackage{fontspec}
\usepackage[utf8]{inputenc}
\usepackage[a4paper,left=.5in,right=.5in,top=.3in,bottom=0.3in]{geometry}
\setlength\parindent{0pt}
\setlength{\parskip}{\baselineskip}
\setmainfont{Helvetica Neue}
\usepackage{hyperref}
\pagestyle{plain}
\begin{document}

\chapter*{Contract language — Version 1 (v1) specification}

Status: draft — this document defines the intended v1 language surface, a
formal grammar (EBNF), lexical rules, the intermediate language (IL) /
bytecode format emitted by the reference compiler, and the VM semantics.
The goal is a stable reference for tooling and a canonical compatibility
target for future implementations.

\section*{Goals for v1}

\begin{itemize}
\item Small, orthogonal feature set suitable for examples and teaching.
\item Clear, deterministic semantics and small predictable bytecode.
\item A simple calling convention to support recursion and local variables.
\item Designed to be easy to parse with a hand-written recursive-descent parser.
\end{itemize}

\section*{Overview}

\begin{itemize}
\item Source files: UTF-8 text, extension \texttt{.ct}.
\item Top-level constructs: Contracts and function declarations.
\item Basic types: integers and strings (no explicit type system in v1).
\item Control flow: if/else, while, switch/case, return.
\item Expressions: literals, identifiers, member access, function calls, binary ops.
\end{itemize}

\section*{Lexical grammar}

Tokens (all whitespace between tokens is ignored except inside strings):

\begin{itemize}
\item IDENTIFIER: [A-Za-z\emph{][A-Za-z0-9}]* (dotted qualified names allowed via '.' between identifiers)
\item INT: [0-9]+ (decimal only)
\item STRING: double quoted, supports escapes like \texttt{\textbackslash{}\textbackslash{}n}, \texttt{\textbackslash{}\textbackslash{}"}, \texttt{\textbackslash{}\textbackslash{}\textbackslash{}\textbackslash{}}.
\item Keywords: \texttt{Contract}, \texttt{fn}, \texttt{switch}, \texttt{case}, \texttt{else}, \texttt{if}, \texttt{while}, \texttt{return}, \texttt{var}
\item Symbols: \texttt{(}, \texttt{)}, \texttt{\{}, \texttt{\}}, \texttt{;}, \texttt{:}, \texttt{,}, \texttt{.}, \texttt{+}, \texttt{-}, \texttt{$<$}, \texttt{$<$=}, \texttt{==}, \texttt{=} (assignment),
\texttt{-$>$} (future use)
\end{itemize}

Comments: \texttt{//} to end-of-line.

All tokens are case-sensitive.

\section*{EBNF grammar (v1)}

```
Start ::= TopLevel*

TopLevel ::= ContractDecl $\vert$ FunctionDecl

ContractDecl ::= 'Contract' IDENTIFIER '\{' TopLevel* '\}'

FunctionDecl ::= 'fn' IDENTIFIER '(' ParamList? ')' Block

ParamList ::= Param (',' Param)*
Param ::= IDENTIFIER (':' IDENTIFIER)?

Block ::= '\{' Statement* '\}'

Statement ::= ExprStatement $\vert$ VarDecl $\vert$ IfStmt $\vert$ WhileStmt $\vert$ SwitchStmt $\vert$ ReturnStmt

ExprStatement ::= Expression ';'
VarDecl ::= 'var' IDENTIFIER (':' IDENTIFIER)? ('=' Expression)? ';'
IfStmt ::= 'if' '(' Expression ')' Block ('else' ':'? (Block $\vert$ Statement))?
WhileStmt ::= 'while' '(' Expression ')' Block
SwitchStmt ::= 'switch' Expression '\{' ( 'case' INT ':' Statement* )* ('else' ':' Statement* )? '\}'
ReturnStmt ::= 'return' Expression? ';'

Expression ::= Assignment
Assignment ::= Equality ( '=' Assignment )?
Equality ::= Relational ( ('==' ) Relational )*
Relational ::= Additive ( ('$<$' $\vert$ '$<$=') Additive )*
Additive ::= Term ( ('+' $\vert$ '-') Term )*
Term ::= Primary
Primary ::= INT $\vert$ STRING $\vert$ IDENTIFIER ('.' IDENTIFIER)* ( '(' ArgList? ')' )?
ArgList ::= Expression (',' Expression)*
``\texttt{
Notes:
- The grammar intentionally keeps operator precedence simple: additive $>$ relational $>$ equality.
- Member access}A.B\texttt{is parsed as a dotted identifier;}A.B(...)` is a call on a qualified name.

\section*{Semantics}

\begin{itemize}
\item No static typing in v1; runtime values are integers or references to strings.
\item Function call semantics: arguments are evaluated left-to-right and pushed on the caller's operand stack. The caller then invokes the callee which maps arguments to its locals in left-to-right order.
\item Variable scoping: functions have locals; \texttt{var} declares a local in the enclosing function or block (implementation-defined in v1 — recommended: function-level locals only).
\item \texttt{switch} uses integer case values; evaluation matches the numeric case and executes the associated statements; \texttt{else} acts as default.
\end{itemize}

\section*{IL / Bytecode format (CIL1)}

File layout (binary):

\begin{itemize}
\item Header: 4 bytes ASCII \texttt{CIL1} + u32 version (1)
\item Constants area: sequence of null-terminated UTF-8 strings
\item Method table: records with method name offset (into constants), arg\emph{count (u16), local}count (u16), code\emph{offset (u32), code}length (u32)
\item Code blob: concatenated bytecode for all methods
\end{itemize}

Instruction set (one-byte opcodes unless followed by immediates):

\begin{itemize}
\item I\_NOP (0x00)
\item I\emph{LDC}I4 (0x01) \texttt{$<$i32 little-endian$>$} — push int32 constant
\item I\emph{LDC}STR (0x02) \texttt{$<$u32 const\_offset$>$} — push referenced constant string address
\item I\emph{PRINT}CONST (0x03) — pop and print const/string
\item I\emph{PRINT}INT (0x04) — pop and print int
\item I\emph{LOAD}LOCAL (0x10) \texttt{$<$u8 index$>$}
\item I\emph{STORE}LOCAL (0x11) \texttt{$<$u8 index$>$}
\item I\emph{ADD (0x20), I}SUB (0x21), I\emph{MUL (0x22), I}DIV (0x23)
\item I\emph{LT (0x30), I}LE (0x31), I\_EQ (0x32)
\item I\_JZ (0x40) \texttt{$<$u32 offset$>$} — jump if top-of-stack zero
\item I\_JMP (0x41) \texttt{$<$u32 offset$>$}
\item I\emph{CALL (0x50) $<$u32 method}index$>$ — call by method table index
\item I\_RET (0x51)
\end{itemize}

Notes:
- Offsets in jump immediates are absolute offsets within the code blob.
- The method table index used by I\_CALL is a 32-bit index into the method table.

\section*{VM semantics}

\begin{itemize}
\item The VM maintains a call stack of frames; each frame contains:

\begin{itemize}
\item return\_ip (offset into code blob),
\item method\_index,
\item locals array (sized per method),
\item operand stack (growable)
\end{itemize}
\item On I\emph{CALL, VM pushes a new frame, copies arguments from the caller's stack into the callee's locals (left-to-right), and transfers control to the callee's code}offset.
\item On I\emph{RET, VM pops the frame and resumes the caller at return}ip. If the language later introduces return values, convention will be to push the result onto the caller's operand stack before returning.
\end{itemize}

Error handling

\begin{itemize}
\item The VM signals runtime errors for:

\begin{itemize}
\item stack underflow/overflow (implementation-defined overflow limit),
\item invalid local index access,
\item invalid instruction,
\item division by zero.
\end{itemize}
\item The reference implementation prints an error and aborts the current program.
\end{itemize}

Versioning and compatibility

\begin{itemize}
\item The header version number is used for format evolution. Tools must refuse to load unknown versions.
\item Language v1 is intentionally conservative: additions are allowed later but removal/change of opcode semantics is forbidden for backward compatibility.
\end{itemize}

Conformance notes (relationship to this repository)

\begin{itemize}
\item The reference C compiler in \texttt{src/} emits the CIL1 format described above.
\item The \texttt{src/runtime.c} implements the VM semantics above.
\item The LSP server uses the reference compiler to emit diagnostics and symbols.
\end{itemize}

Examples

\begin{itemize}
\item A simple factorial function in Contract syntax and the approximate bytecode emitted by the reference compiler.
\end{itemize}

Contract source:

\texttt{ct
Contract Math \{
  fn fact(n) \{
    var acc = 1;
    while (n $>$ 1) \{
      acc = acc * n;
      n = n - 1;
    \}
    return acc;
  \}
\}
}

Corresponding (conceptual) IL:

\begin{itemize}
\item method table: Math.fact arg\emph{count=1 local}count=2 code\emph{offset=0 code}length=...
\item code at offset 0:
I\emph{LOAD}LOCAL 0    ; n
I\emph{LDC}I4 1
I\emph{LT
I}JZ \texttt{$<$after\_loop$>$}
I\emph{LOAD}LOCAL 1    ; acc
I\emph{LOAD}LOCAL 0
I\emph{MUL
I}STORE\emph{LOCAL 1
I}LOAD\emph{LOCAL 0
I}LDC\emph{I4 1
I}SUB
I\emph{STORE}LOCAL 0
I\emph{JMP $<$loop}top$>$
\texttt{$<$after\_loop$>$}:
I\emph{LOAD}LOCAL 1
I\_RET
\end{itemize}

Future extensions (non-normative)

\begin{itemize}
\item Add a small type system with optional annotations.
\item Add first-class strings and string ops in the IL.
\item Add optimized opcodes for small integer constants (I\emph{LDC}0..I\emph{LDC}7).
\end{itemize}

Feedback and process

This spec is intended as the canonical v1. If you'd like, I can:

\begin{itemize}
\item convert the EBNF to an ANTLR grammar or a PEG for automated parsing,
\item generate a test-suite of small example programs and expected diagnostics/IL,
\item expand the IL to include debug info (line/col mapping) for improved tooling.
\end{itemize}

\end{document}
